\documentclass[12pt]{article}

\usepackage{hyperref}
\usepackage{graphicx}
\usepackage{booktabs}
\usepackage{geometry}
\usepackage{tabularx}
\usepackage{float}
\usepackage[none]{hyphenat}
\usepackage{setspace}
\usepackage{subcaption}
\usepackage{titlesec}
\usepackage{amsmath}

\geometry{margin=1in}
\setstretch{0.96}

\titlespacing*{\section}{0pt}{10pt}{4pt}

\setlength{\parindent}{0pt}
\setlength{\parskip}{6pt}

\title{Labwork 3}
\author{Trinh Thanh Vinh - 2540110}
\date{\today}

\begin{document}

\maketitle
\section{Introduction}

This report presents the implementation of a simple neural network using the backpropagation algorithm in Python. The neural network was implemented from scratch without using libraries such as TensorFlow or PyTorch. The objective of this experiment is to train the neural network to solve the XOR classification problem.

\section{Neural Network Structure}

The neural network consists of input layer, hidden layer, and output layer. The network structure is loaded from a file named \texttt{struc.txt}. For the XOR experiment, the structure contains 3 layers:
\begin{itemize}
    \item 2 input neurons
    \item 3 hidden neurons
    \item 1 output neuron
\end{itemize}

\section{Feedforward Process}

Each neuron computes:

\[
z = \sum_{i=1}^{n} w_i x_i + b
\]

where:
\begin{itemize}
    \item $w_i$ are the weights
    \item $x_i$ are the inputs
    \item $b$ is the bias
\end{itemize}

The sigmoid activation function is used:

\[
\sigma(x) = \frac{1}{1 + e^{-x}}
\]

The neuron output becomes:

\[
a = \sigma(z)
\]

\section{Backpropagation Implementation}

Backpropagation is used to compute gradients for updating weights and biases.

\subsection{Output Layer Delta}

For the output neuron, the delta is computed as:

\[
\delta = \hat{y} - y
\]

where:
\begin{itemize}
    \item $\hat{y}$ is the predicted output
    \item $y$ is the true label
\end{itemize}

In the implementation:

\begin{verbatim}
s = neuron.sigmoid(-neuron.last_z)
neuron.delta = -y_true[j] + (1 - s)
\end{verbatim}

Since:

\[
1 - \sigma(-z) = \sigma(z)
\]

the equation becomes:

\[
\delta = \sigma(z) - y
\]

\subsection{Hidden Layer Delta}

For hidden neurons:

\[
\delta_j = \sigma(z_j)(1 - \sigma(z_j))
\sum_k w_{jk}\delta_k
\]

where:
\begin{itemize}
    \item $w_{jk}$ is the weight connecting hidden neuron $j$ to neuron $k$
    \item $\delta_k$ is the delta from the next layer
\end{itemize}

Implementation:

\begin{verbatim}
grad_sum += next_neuron.weights[j] * next_neuron.delta
neuron.delta = (1 - s) * s * grad_sum
\end{verbatim}

\subsection{Gradient Descent Update}

Weights and biases are updated using gradient descent:

\[
w = w - \eta \delta x
\]

\[
b = b - \eta \delta
\]

where:
\begin{itemize}
    \item $\eta$ is the learning rate
\end{itemize}

Implementation:

\begin{verbatim}
neuron.weights[k] =
    neuron.weights[k]
    - l * neuron.delta * neuron.last_input[k]

neuron.bias =
    neuron.bias - l * neuron.delta
\end{verbatim}

\section{Loss Function}

The binary cross-entropy loss function is used:

\[
L = -\frac{1}{N}
\sum_{i=1}^{N}
\left(
y_i z_i - \log(1 + e^{z_i})
\right)
\]

where:
\begin{itemize}
    \item $N$ is the number of samples
    \item $z_i$ is the neuron output before activation
\end{itemize}

\section{XOR Dataset Experiment}

\subsection{Dataset}

The XOR dataset:

\[
(0,0) \rightarrow 0
\]

\[
(0,1) \rightarrow 1
\]

\[
(1,0) \rightarrow 1
\]

\[
(1,1) \rightarrow 0
\]

\subsection{Training Configuration}

\begin{itemize}
    \item Learning rate: 0.5
    \item Maximum epochs: 50000
    \item Activation function: Sigmoid
\end{itemize}

\section{Results}

After 50000 epochs, the network produced the following predictions:

\begin{verbatim}
[0, 0] -> predicted: 8.051108211763272e-05
expected: 0

[0, 1] -> predicted: 0.999916280093545
expected: 1

[1, 0] -> predicted: 0.9999180963066683
expected: 1

[1, 1] -> predicted: 0.0001239938381239292
expected: 0
\end{verbatim}

\section{Analysis}

The neural network successfully learned the XOR function, with the predictions being extremely close to the expected outputs. This demonstrates that the backpropagation algorithm was implemented correctly, allowing the network to adjust its weights and biases to minimize the loss function effectively. The hidden layer allows the network to capture the non-linear relationships in the XOR dataset, which is not linearly separable. The gradient descent successfully minimized the loss function, leading to accurate predictions after training.
\end{document}